% Pursuing Stacks (cotes 134-2 to 134-8): archive cote 134-6, site page 73 (PDF page 74), typescript page 391.
% typescript. Draft compared: ps5.tex:5103-5189.
% AI-transcribed by Claude Opus 5.5 (claude-opus-5-5, adaptive thinking, effort high), Message Batches API, batch ps3 (msgbatch_01XdFmuVVBAcgc7EFxdtCqsj), 2026-11-12; reviewed against both supplied scan halves; see CHECK_DECISIONS.tsv.
\scanpage{134-6/73}
\mspage{391}

which will be \xout{handy} convenient mainly in the context of the next subsecti-%
on (when \hand{$\mathcal{A}$} is of the type A\^{}).

T\(\overset{\text{\tiny -\,-}}{\underset{\text{h}}{\text{h}}}\)e exactness property needed for the functor (21) is equiva lent
with the property that for any object H in \hand{$\mathcal{A}$}$_{\mathrm{ab}}$ , the functor
\[ G \hand{\longmapsto} \mathrm{Hom}(F\otimes G,\ H) \hand{\simeq} \mathrm{Bil}(F,G;\ H) \]
from \hand{$\mathcal{A}$}$_{\mathrm{ab}}{}^{\mathrm{o}}$ to (Sets) commutewith small inverse limits. As \hand{$\mathcal{A}$}$_{\mathrm{ab}}$ is a
pseudotopos \inserted[typed]{(prop.2)}, this is equivalent \xout{with} by prop.1 with this functor being
representable, \over{B}{B}\over{y}{y} definition of Bil, \inserted[typed]{we get} \xout{thisx\allowbreak{}functorx\allowbreak{}isx\allowbreak{}justx\allowbreak{}thex\allowbreak{}functor}
\xout{\hand{\uline{\textcolor{black}{H\over{o}{o}m}}}x\allowbreak{}x\allowbreak{}$_{\mathrm{ab}}$\hand{(}F\hand{,}H\hand{)}} \[ \mathrm{Bil}(F,G;H)\hand{\simeq} \mathrm{Hom}_{\hand{\mathcal{A}}^{\wedge}_{\mathrm{ab}}}(\over{G}{G},\ \underline{\mathrm{Hom}}_{\underline{Z}}(F,H)) \ , \]
where the object
\[ (24) \qquad \underline{\mathrm{Hom}}_{\underline{Z}}(F,H) \]
is taken in \xout{the} \inserted{the category of presh\over{\illegible{1}}{e}aves} \hand{$\mathcal{A}$}$^{\wedge}_{\mathrm{ab}}$ (cheating a little with univer\over[typed]{e}{s}\over[typed]{s}{e}s here ...). To sum
up, the condition we want for (21) just boils down to the representabili-%
ty of the \inserted[typed]{abelian} group objects (24) in \hand{$\mathcal{A}$}$^{\wedge}_{\mathrm{ab}}$ , for any two objects F,H in \hand{$\mathcal{A}$}$_{\mathrm{ab}}$ ,
i.e. essentially to the existence of ``internal \uline{Hom}'s '' in the category
\hand{$\mathcal{A}$}$_{\mathrm{ab}}$\hand{(}endowed with the tensor product $\over{\otimes}{\otimes}_{\underline{Z}}$ \hand{)}, satisfying the familiar
Cartan isomorphism formula
\[ (25) \qquad \mathrm{Hom}(F\otimes G,\ H)\hand{\simeq} \mathrm{Hom}(G,\ \underline{\mathrm{Hom}}(F,H)) \ . \]

To sum up, what is needed for a nice formalism of interior and
exterior tensor products and Hom's for \xout{sheaves} ``sheaves'' and ``presheaves''
on the pseudotopos \hand{$\mathcal{A}$} , \over{a}{a}\over{r}{r}e the following assumptions\deleted{:} \inserted{on $\mathcal{A}$} :

1) T\over[typed]{n}{e}\over[typed]{e}{n}sor products and \inserted[typed]{corresponding} internal \uline{Hom}'s exist in \hand{$\mathcal{A}$}$_{\mathrm{ab}}$ , and

2) the abelianization functor Wh$_{\hand{\mathcal{A}}}$ (15) exists ,

the latter assumption being needed in order to feel at ease with the
equivalence between the categories (19) of sheaves and cosheaves, and
their ``abelianized'' interpretations (19').

Under these assumptions, we define exterior tensor product opera-%
tion of \hand{$\mathcal{A}$}$_{\mathrm{ab}}$ upon a category of \over[typed]{\illegible{1}}{\unsure{c|s}}osheaves \hand{$\mathcal{A}^{\mathcal{M}}$}
\margin{$\circledast$} \[ (26) \qquad \hand{\mathcal{A}}_{\mathrm{ab}} \times \hand{\mathcal{A}^{\mathcal{M}}} \hand{\longrightarrow} \hand{\mathcal{A}^{\mathcal{M}}}, \qquad (F,L)\hand{\mapsto} F \hand{\mathbin{\tikz[baseline=-.5ex]{\node[circle,draw,inner sep=.04em] {\ensuremath{\textcolor{black}{*}}};}}} L \quad , \]
by formu\over[typed]{b}{l}a (22) (which may be written also under the form (22')). This
operation has the obvious associativity property
\[ (27) \qquad F \hand{\mathbin{\tikz[baseline=-.5ex]{\node[circle,draw,inner sep=.04em] {\ensuremath{\textcolor{black}{*}}};}}} (G \hand{\mathbin{\tikz[baseline=-.5ex]{\node[circle,draw,inner sep=.04em] {\ensuremath{\textcolor{black}{*}}};}}} L) \hand{\simeq} (F \over{\otimes}{\otimes} G) \hand{\mathbin{\tikz[baseline=-.5ex]{\node[circle,draw,inner sep=.04em] {\ensuremath{\textcolor{black}{*}}};}}} L \quad , \]
and moreover the unit $\underline{Z}_{\hand{\mathcal{A}}}$ for the internal tensor product in \hand{$\mathcal{A}$}$_{\mathrm{ab}}$ operates
on \hand{$\mathcal{A}^{\mathcal{M}}$} as the identity functor, i.e.

\margin{73}
